\noindent\textbf{Extended Data Fig.~X | Reconstruction uncertainty in the association between phenomic diversity and diversification.}
\textbf{a,b}, Distributions of Pearson correlation coefficients (\textbf{a}) and two-sided nominal $p$-values (\textbf{b}) between undetrended $N_{\mathrm{eff}}^{+}$ and family-level PyRate net diversification trajectories.
\textbf{c,d}, Corresponding distributions for BAMM after linear detrending.
For the BAMM comparison, each $N_{\mathrm{eff}}^{+}$ and BAMM trajectory was independently regressed against time-bin midpoint age, and correlations were calculated between the residuals, equivalent to partial Pearson correlations controlling for linear time.
Net diversification was calculated as $\lambda-\mu$ within each diversification sample.
Each correlation used 13 complete 20-million-year bins spanning 260--0~Ma.
All 1,000 $N_{\mathrm{eff}}^{+}$ stochastic-map reconstructions were paired with all 2,000 retained PyRate trajectories or all 18,001 retained BAMM trajectories, yielding 2,000,000 and 18,001,000 trajectory pairs, respectively.
PyRate trajectories comprised the final 200 saved samples from each of 10 runs; BAMM trajectories were reconstructed from the existing posterior on the full 765-tip phylogeny.
$N_{\mathrm{eff}}^{+}$ was calculated on the 77-taxon phenomic tree.
Each pairing retained the complete sequence of 13 bins.
Histograms show percentages of trajectory pairs; $p$-values are displayed on a logarithmic scale.
Dashed grey lines indicate correlations or tests of the mean trajectories, with the same detrending applied in the BAMM panels; dotted orange lines mark $p=0.05$.
Annotations show median correlations and central 95\% reconstruction ranges, or percentages of pairs with $p<0.05$.
These ranges describe variation among reconstructions, rather than sampling confidence intervals.
The two-sided tests used 11 degrees of freedom for PyRate and 10 for BAMM, accounting for the time covariate in the latter.
Trajectory pairs are not independent experimental replicates.
The tests remain nominal: linear detrending does not guarantee removal of temporal autocorrelation.
